\section{財政状況報告}

\begin{fillblank}[自作 2026/11/7]
  「公的年金財政状況報告 -令和6(2024)年度- ポイント」に記載の文章の以下の空欄を埋めなさい。
\\

  1 公的年金の収支状況

  公的年金制度全体でみると、令和6(2024)年度は、運用損益分 を除いた収入総額 55.7 兆円、
  支出総額 55.8 兆円であったこと から、運用損益分を除いた単年度収支残は 0.1 兆円のマイナス。 
  また、時価ベースの運用損益は 2.0 兆円のプラス。

  その結果、時価ベースの年度末積立金は前年度に比べ 2.0 兆円 増加し、306.0 兆円。
\\

2 公的年金の財政状況の評価

令和6(2024)年度までの実績と令和6(2024)年財政検証の前提や将来
見通しを比較するだけではなく、長期的な財政の均衡の観点から評価。

\begin{itemize}
  \item 令和6(2024)年度は令和6(2024)年財政検証における将来見通しと比較する初年度であるが、
  本報告では令和6(2024)年 度の実績について以下のことを確認した。
  \begin{enumerate}[(1)]
    \item 国民年金第1号被保険者数、厚生年金被保険者数ともに実績 が将来見通しを上回った。
    これは、（　　A　　）の実績 が令和5(2023)年推計におけるいずれの仮定値も上回って
    いることも影響していると考えられる。また、65 歳の平均余 命の実績は、
    令和5(2023)年推計における死亡（　　B　　）の仮定値 と概ね同水準であった。
    これらは公的年金財政にはプラスの 効果となる。

    \item 一方で、合計特殊出生率は平成 28(2016)年より低下傾向が 続いており、
    令和6(2024)年の実績は令和5(2023)年推計に おける出生（　　C　　）の仮定値と概ね同水準となっている。
    また、 （　　D　　）は令和6(2024)年財政検証におけるいずれのケースの前提も下回っている。
    これらは公的年金財政にはマイナスの効果となる。

    \item 積立金については、令和5(2023)年度末の実績が将来見通しを上回っていたこともあり、
    令和6(2024)年度末においても実績が将来見通しを上回った。
    なお、令和5(2023)年度末の実績が将来見通しを上回った部分のほとんどは、
    将来見通しの積立金に平滑化したものが使われていることによるものである。

    \item 令和6(2024)年度は、マクロ経済スライドによる給付水準調整が行われた一方で、
    賃金の上昇が物価上昇に追いつかず、 実質賃金の伸びはマイナスとなったことから、
    令和5(2023) 年度とは異なり、既裁定年金の伸びを賃金の伸びより抑制する効果は発動されなかった。
  \end{enumerate}

  \item これらの将来見通しからの乖離が、一時的なものではなく中長期的に続いた場合には、
  年金財政に与える影響は大きなものと なる。
  特に、令和6(2024)年の合計特殊出生率は令和5(2023) 年より低下し、
  令和5(2023)年推計における出生低位の仮定値 と概ね同水準である。
  \\ このような傾向が今後も継続するようで あれば、将来の年金制度の運営は大きな影響を受ける。

  \item 年金財政の観点からは、人口要素、経済要素等いずれも短期的な動向にとらわれることなく、長期的な観点から財政状況の動 向を注視すべきである
\end{itemize}


\end{fillblank}

\begin{answerbox}
  \begin{enumerate}[A: ]
    \item 外国人入国超過数
    \item 高位
    \item 低位
    \item 実質賃金上昇率（対物価）
  \end{enumerate}
\end{answerbox}

\begin{explanationbox}
  ＜参考＞公的年金財政状況報告

  \url{https://www.mhlw.go.jp/stf/seisakunitsuite/bunya/0000128082_00001.html}
\end{explanationbox}

\newpage